% Контрольное восстановление §1.16 предоставленного учебника И. В. Яковлева.
% Текст и все рисунки редактируемы. Основной файл для LaTeX Workshop.
\documentclass[12pt,a4paper]{article}
\usepackage[T2A]{fontenc}
\renewcommand{\rmdefault}{cmr}
\renewcommand{\familydefault}{\rmdefault}
\usepackage{newunicodechar}
\newunicodechar{А}{\char"C0 }
\newunicodechar{Б}{\char"C1 }
\newunicodechar{В}{\char"C2 }
\newunicodechar{Г}{\char"C3 }
\newunicodechar{Д}{\char"C4 }
\newunicodechar{Е}{\char"C5 }
\newunicodechar{Ж}{\char"C6 }
\newunicodechar{З}{\char"C7 }
\newunicodechar{И}{\char"C8 }
\newunicodechar{Й}{\char"C9 }
\newunicodechar{К}{\char"CA }
\newunicodechar{Л}{\char"CB }
\newunicodechar{М}{\char"CC }
\newunicodechar{Н}{\char"CD }
\newunicodechar{О}{\char"CE }
\newunicodechar{П}{\char"CF }
\newunicodechar{Р}{\char"D0 }
\newunicodechar{С}{\char"D1 }
\newunicodechar{Т}{\char"D2 }
\newunicodechar{У}{\char"D3 }
\newunicodechar{Ф}{\char"D4 }
\newunicodechar{Х}{\char"D5 }
\newunicodechar{Ц}{\char"D6 }
\newunicodechar{Ч}{\char"D7 }
\newunicodechar{Ш}{\char"D8 }
\newunicodechar{Щ}{\char"D9 }
\newunicodechar{Ъ}{\char"DA }
\newunicodechar{Ы}{\char"DB }
\newunicodechar{Ь}{\char"DC }
\newunicodechar{Э}{\char"DD }
\newunicodechar{Ю}{\char"DE }
\newunicodechar{Я}{\char"DF }
\newunicodechar{а}{\char"E0 }
\newunicodechar{б}{\char"E1 }
\newunicodechar{в}{\char"E2 }
\newunicodechar{г}{\char"E3 }
\newunicodechar{д}{\char"E4 }
\newunicodechar{е}{\char"E5 }
\newunicodechar{ж}{\char"E6 }
\newunicodechar{з}{\char"E7 }
\newunicodechar{и}{\char"E8 }
\newunicodechar{й}{\char"E9 }
\newunicodechar{к}{\char"EA }
\newunicodechar{л}{\char"EB }
\newunicodechar{м}{\char"EC }
\newunicodechar{н}{\char"ED }
\newunicodechar{о}{\char"EE }
\newunicodechar{п}{\char"EF }
\newunicodechar{р}{\char"F0 }
\newunicodechar{с}{\char"F1 }
\newunicodechar{т}{\char"F2 }
\newunicodechar{у}{\char"F3 }
\newunicodechar{ф}{\char"F4 }
\newunicodechar{х}{\char"F5 }
\newunicodechar{ц}{\char"F6 }
\newunicodechar{ч}{\char"F7 }
\newunicodechar{ш}{\char"F8 }
\newunicodechar{щ}{\char"F9 }
\newunicodechar{ъ}{\char"FA }
\newunicodechar{ы}{\char"FB }
\newunicodechar{ь}{\char"FC }
\newunicodechar{э}{\char"FD }
\newunicodechar{ю}{\char"FE }
\newunicodechar{я}{\char"FF }
\newunicodechar{Ё}{\char"9C }
\newunicodechar{ё}{\char"BC }
\newunicodechar{—}{\textemdash}
\newunicodechar{–}{\textendash}
\newunicodechar{«}{\guillemotleft}
\newunicodechar{»}{\guillemotright}
\newunicodechar{…}{\ldots}

\DeclareFontShape{T2A}{cmr}{m}{n}{<8>larm0800 <10.95>larm1095 <12>larm1200 <14.4>larm1200 <17.28>larm1728}{}
\DeclareFontShape{T2A}{cmr}{bx}{n}{<14.4>labx1440 <17.28>labx1728 <12>labx1200}{}
\usepackage[russian]{babel}
\usepackage{amsmath,amssymb}
\usepackage[left=43.654bp,textwidth=510.236bp,top=52.169bp,textheight=750bp,footskip=7.161bp]{geometry}
\usepackage{xcolor,tikz,wrapfig}
\definecolor{booklink}{rgb}{0,.3,.9}
\definecolor{bookgray}{rgb}{.9,.9,.9}
\usepackage[unicode,hypertexnames=false,colorlinks=true,linkcolor=blue,urlcolor=blue]{hyperref}
\AtBeginDvi{\special{pdf:mapfile +cm-super-t2a.map}}
\setlength{\parindent}{17.625pt}
\setlength{\parskip}{0pt}
\setlength{\emergencystretch}{0pt}
\pagestyle{plain}
% Native vector geometry measured from the reference in big points.
% Ground-left origin, y upwards. Source heads are filled narrow triangles.
\newcommand{\WorkParallelDiagram}{%
\begin{tikzpicture}[x=1bp,y=1bp,line width=.5bp,
 every node/.style={inner sep=0pt,outer sep=0pt,font=\fontsize{10}{12}\selectfont}]
 \path[use as bounding box] (-.25,-30.25) rectangle (212.511,21);
 \foreach \x in {30,170} {
   \fill[bookgray] (\x-15,0) rectangle (\x+15,20);
   \draw (\x-15,0)--(\x-15,20)--(\x+15,20)--(\x+15,0);
 }
 \draw (0,0)--(200,0);
 \foreach \x in {0,3,...,198} {\draw (\x,0)--++(-3,-3);}
 \foreach \x in {30,170} {
   \draw[dash pattern=on 1.5bp off 1.5bp] (\x,10)--(\x,-30);
 }
 \draw[blue] (30,-26)--(170,-26);
 \filldraw[blue] (163.106,-24.7845)--(170,-26)--(163.106,-27.2155)--cycle;
 \node[anchor=base west] at (97.6646,-23) {$\vec s$};
 \foreach \x in {30,170} {
   \fill[red] (\x,10) circle[radius=1.5bp];
   \draw[red] (\x,10)--(\x+40,10);
   \filldraw[red] (\x+33.106,11.2155)--(\x+40,10)--(\x+33.106,8.7845)--cycle;
   \node[anchor=base west] at (\x+36.105,13) {$\vec F$};
 }
\end{tikzpicture}}

\newcommand{\WorkObliqueDiagram}{%
\begin{tikzpicture}[x=1bp,y=1bp,line width=.5bp,
 every node/.style={inner sep=0pt,outer sep=0pt,font=\fontsize{10}{12}\selectfont}]
 \path[use as bounding box] (-.25,-30.25) rectangle (268,62);
 \foreach \x in {50,190} {
   \fill[bookgray] (\x-20,0) rectangle (\x+20,20);
   \draw (\x-20,0)--(\x-20,20)--(\x+20,20)--(\x+20,0);
 }
 \draw (0,0)--(240,0);
 \foreach \x in {0,3,...,240} {\draw (\x,0)--++(-3,-3);}
 \foreach \x in {50,190} {
   \draw[dash pattern=on 1.5bp off 1.5bp] (\x,10)--(\x,-30);
 }
 \draw[blue] (50,-26)--(190,-26);
 \filldraw[blue] (183.106,-24.7845)--(190,-26)--(183.106,-27.2155)--cycle;
 \node[anchor=base west] at (117.6646,-23) {$\vec s$};
 \foreach \x in {50,190} {
   \fill[red] (\x,10) circle[radius=1.5bp];
   \draw[red] (\x,10)--(\x+60,10);
   \filldraw[red] (\x+53.106,11.2155)--(\x+60,10)--(\x+53.106,8.7845)--cycle;
   \draw[red] (\x,10)--(\x,50);
   \filldraw[red] (\x-1.2155,43.106)--(\x,50)--(\x+1.2155,43.106)--cycle;
   \draw[red] (\x,10)--(\x+60,50);
   \filldraw[red] (\x+54.9382,45.16455)--(\x+60,50)--(\x+53.58966,47.18735)--cycle;
   \draw[dash pattern=on 1.5bp off 1.5bp] (\x+60,50)--(\x+60,10);
   \draw[dash pattern=on 1.5bp off 1.5bp] (\x+60,50)--(\x,50);
   \draw (\x+30,10) arc[start angle=0,end angle=33.6900675,radius=30bp];
   \node[anchor=base west] at (\x+63,6.953) {$\vec F_{\parallel}$};
   \node[anchor=base west] at (\x-16.132,45.929) {$\vec F_{\perp}$};
   \node[anchor=base west] at (\x+62.1,52.1) {$\vec F$};
   \node[anchor=base west] at (\x+31.712,16.549) {$\alpha$};
 }
\end{tikzpicture}}

% Native TikZ figure 1.50; geometry measured from reference PDF page 83.
% Ground-left source coordinate (408.495,678.173) bp, y upwards here.
\newcommand{\ReplicaGravityDiagram}{%
\begin{tikzpicture}[x=1bp,y=1bp,line width=.5bp,
 every node/.style={inner sep=0pt,outer sep=0pt,font=\fontsize{10}{12}\selectfont}]
  \coordinate (P) at (90,130);
  \coordinate (Q) at (10,40);
  \draw (0,0)--(100,0);
  \foreach \x in {0,3,...,99} {\draw (\x,0)--++(-3,-3);}
  \fill[blue] (P) circle[radius=1.5bp];
  \draw[blue] (P)--(Q);
  \filldraw[blue] (13.67166,45.96033)--(Q)--(15.48874,44.34509)--cycle;
  \draw[dash pattern=on 1.5bp off 1.5bp] (P)--(90,0);
  \draw[dash pattern=on 1.5bp off 1.5bp] (Q)--(10,0);
  \draw[dash pattern=on 1.5bp off 1.5bp] (Q)--(90,40);
  \node[anchor=base west] at (93,126.596) {$P$};
  \node[anchor=base west] at (-.876,37.564) {$Q$};
  \node[anchor=base west] at (43.2296,87.100) {$\vec s$};
  \node[anchor=base west] at (93,62.287) {$h_1$};
  \node[anchor=base west] at (13,17.287) {$h_2$};
  \draw (90,110) arc[start angle=270,end angle=228.37,radius=20];
  \node[anchor=base west] at (78.687,104.016) {$\alpha$};
  \draw[red] (10,120)--(10,90);
  \filldraw[red] (8.7845,96.89374)--(10,90)--(11.2155,96.89374)--cycle;
  \node[anchor=base west] at (13.00238,102.409) {$\vec g$};
  % Preserve the old macro's TeX box and baseline, so text flow stays fixed.
  \pgfresetboundingbox
  \pgfpathrectangle{\pgfpoint{-3.2bp}{-3.2bp}}{\pgfpoint{114.85211pt}{141.11418pt}}
  \pgfusepath{use as bounding box}
\end{tikzpicture}%
}

% Калибровка высоты двух иллюстраций по строкам оригинала (единица bp).
% Рисунки остаются собственными TikZ-объектами; здесь только размещение.
\newcommand{\ReplicaPlacedWorkFigure}[5]{%
  \par\begingroup
  \setbox0=\hbox{#1}%
  \dimen0=\prevdepth
  \kern\dimexpr #2bp-\dimen0-\ht0+30.25bp\relax
  \nointerlineskip
  \hbox to\textwidth{\hspace*{#3bp}\box0\hfil}%
  \setbox1=\hbox{#4}%
  \kern\dimexpr #5bp-\ht1\relax
  \nointerlineskip\hbox to\textwidth{\hfil\box1\hfil}%
  \endgroup
}
\newcommand{\ReplicaWorkFigureOne}{%
  \ReplicaPlacedWorkFigure{\WorkParallelDiagram}{45.256}{149.118}{Рис.~1.48. $A=Fs$}{24.409}%
  \vskip10.384bp\relax
}
\newcommand{\ReplicaWorkFigureTwo}{%
  \ReplicaPlacedWorkFigure{\WorkObliqueDiagram}{84.256}{124.118}{Рис.~1.49. $A=Fs\cos\alpha$}{24.409}%
  \vskip11.466bp\relax
}

\begin{document}
\setcounter{page}{80}
\section*{1.16\quad Энергия}
Мы приступаем к изучению энер\-гии — фундаментального физического понятия. Но предвари\-тельно нужно разобраться с другой физической величиной — работой силы.

\kern-.149bp\relax
\subsection*{1.16.1\quad Работа}
Пусть на тело действует постоянная сила $\vec F$ и тело, двигаясь прямолинейно по горизонтальной поверхности, совершило перемещение $\vec s$. Сила $\vec F$ не обязательно является непосредственной при\-чиной перемещения (так, сила тяже\-сти не является непосредственной при\-чиной перемещения шкафа, который передвигают по комнате).

Предположим сначала, что векторы силы и перемещения сонаправлены (рис.~\textcolor{booklink}{1.48}; остальные силы, действующие на тело, не указаны).
\ReplicaWorkFigureOne

В этом простейшем случае работа $A$ определяется как произведение модуля силы на модуль перемещения:
\begin{equation}A=Fs.\tag{1.76}\end{equation}
\vskip-.432bp\relax


Единицей измерения работы служит \textit{джоуль} (Дж): Дж = Н $\cdot$ м. Таким образом, если под действием силы 1 Н тело перемещается на 1 м, то сила совершает работу 1 Дж.

Работа силы, перпендикулярной перемещению, по определению считается равной нулю. Так, в данном случае сила тяже\-сти и сила реакции опоры не соверша\-ют работы.

Пусть теперь вектор силы образует с вектором перемещения острый угол $\alpha$ (рис.~\textcolor{booklink}{1.49}).
\ReplicaWorkFigureTwo

Разложим силу $\vec F$ на две составляющие: $\vec F_{\parallel}$ (параллельную перемещению) и $\vec F_{\perp}$ (перпенди\-кулярную перемещению). Работу совершает только $\vec F_{\parallel}$. Поэтому для работы силы $\vec F$ получаем: $A=F_{\parallel}s=F\cos\alpha\cdot s$. Итак,
\begin{equation}A=Fs\cos\alpha.\tag{1.77}\end{equation}
\vskip-.432bp\relax


Если вектор силы образует с вектором перемещения тупой угол $\alpha$, то работа по-прежнему определяется формулой (\textcolor{booklink}{1.77}). В этом случае работа оказывается отрицательной.

Например, работа силы трения скольжения, действующей на тело в рассмотренных ситуа\-циях, будет отрицательной, так как сила трения направлена противоположно перемещению. В этом случае имеем: $\alpha=180^{\circ}$, $\cos\alpha=-1$, и для работы силы трения получаем:
\vadjust{\kern-1.007bp}
$$A_{\text{тр}}=-F_{\text{тр}}s=-\mu mgs,$$
\newpage
\noindent где $m$ — масса тела, $\mu$ — коэффициент трения между телом и опорой.

\vspace{5.146bp}
Соотношение (\textcolor{booklink}{1.77}) означает, что работа является \textit{скалярным произведением} векторов силы и перемещения:
$$A=\vec F\vec s.$$

Это позволяет вычислять работу через координаты данных векторов:
$$A=F_xs_x+F_ys_y+F_zs_z.$$

Пусть на тело действуют несколько сил $\vec F_1,\vec F_2,\ldots,\vec F_n$, и $\vec F$ — равнодействующая этих сил. Для работы силы $\vec F$ имеем:
$$A=\vec F\vec s=(\vec F_1+\vec F_2+\ldots+\vec F_n)\vec s=\vec F_1\vec s+\vec F_2\vec s+\ldots+\vec F_n\vec s,$$
\noindent или
$$A=A_1+A_2+\ldots+A_n,$$
\noindent где $A_1,A_2,\ldots,A_n$ — работы сил $\vec F_1,\vec F_2,\ldots,\vec F_n$. Итак, \textit{работа равнодействующей приложенных к телу сил равна сумме работ каждой силы в отдельности.}

\subsection*{1.16.2\quad Мощность}
Часто имеет значение быстрота, с ко\-торой совершается работа. Скажем, на практике важно знать, какую работу сможет выполнить данное устройство за фиксированное время.

\textit{Мощность} — это величина, характеризующая скорость совершения работы. Мощность $N$ есть отношение работы $A$ ко времени $t$, за которое эта работа совершена:
$$N=\frac{A}{t}.$$

Мощность измеряется в \textit{ваттах} (Вт). 1 Вт = 1 Дж/с, то есть 1 Вт — это такая мощность, при ко\-торой работа в 1 Дж совершается за 1 с.

Предположим, что силы, действующие на тело, уравновешены, и тело движется равномерно и прямолинейно со скоростью $\vec v$. В этом случае существует полезная формула для мощности, развиваемой одной из действующих сил $\vec F$.

За время $t$ тело совершит перемещение $\vec s=\vec vt$. Работа силы $\vec F$ будет равна:
$$A=\vec F\vec s=\vec F\vec vt.$$
\noindent Отсюда получаем мощность:
$$N=\vec F\vec v,$$
\noindent или
$$N=Fv\cos\alpha,$$
\noindent где $\alpha$ — угол между векторами силы и скорости.

Наиболее часто эта формула используется в ситуации, когда $\vec F$ — «сила тяги» двигателя автомобиля (которая на самом деле есть сила трения ведущих колёс о дорогу). В этом случае $\alpha=0$, и мы получаем просто:
$$N=Fv.$$
\newpage
\subsection*{1.16.3\quad Механическая энергия}
Энергия является мерой движе\-ния и взаимодействия любых объектов в природе. Имеются различные формы энер\-гии: механиче\-ская, тепловая, электромагнитная, ядерная\ldots

Опыт показывает, что энергия не появляется ниоткуда и не исчезает бесследно, она лишь переходит из одной формы в другую. Это самая общая формулировка \textit{закона со\-хране\-ния энер\-гии.}

Каждый вид энер\-гии представляет собой некоторое математическое выражение. Закон со\-хране\-ния энер\-гии означает, что в каждом явлении природы определённая сумма таких выра\-жений остаётся постоянной с течением времени.

Измеряется энергия в джоулях, как и работа.

\textit{Механическая энергия} является мерой движе\-ния и взаимодействия механических объектов (материальных точек, твёрдых тел).

Мерой движе\-ния тела является \textit{кинетическая энергия}. Она зависит от скорости тела. Мерой взаимодействия тел является \textit{потенциаль\-ная энергия}. Она зависит от взаимного расположения тел.

Механическая энергия си\-стемы тел равна сумме кинетической энер\-гии тел и потенциаль\-ной энер\-гии их взаимодействия друг с другом.

\kern-.487bp\relax
\subsection*{1.16.4\quad Кинетическая энергия}
Кинетической энергией тела (принимаемого за материальную точку) называется величина
\vskip-2.801bp\relax
$$K=\frac{mv^2}{2}\,,$$
\vskip2.181bp\relax
\noindent где $m$ — масса тела, $v$ — его скорость.

Кинетической энергией си\-стемы из $N$ тел называется сумма кинетических энергий каждого тела:
$$K=\frac{m_1v_1^2}{2}+\frac{m_2v_2^2}{2}+\ldots+\frac{m_Nv_N^2}{2}\,.$$
\vskip-1.413bp\relax

Если тело движется под действием силы $\vec F$, то кинетическая энергия тела, вообще гово\-ря, ме\-няется со временем. Оказывается, именение кинетической энер\-гии тела за некоторый промежуток времени равно работе силы $\vec F$. Покажем это для случая прямолинейного равно\-ускоренного движе\-ния.

Пусть $\vec v_1$ — начальная скорость, $\vec v_2$ — конечная скорость тела. Выберем ось $X$ вдоль тра\-ек\-тории тела (и, соответственно, вдоль вектора силы $\vec F$). Для работы силы $\vec F$ получаем:
\vskip-2.801bp\relax
$$A=\vec F\vec s=F_xs_x=ma_xs_x=ma_x\frac{v_{2x}^2-v_{1x}^2}{2a_x}=\frac{mv_{2x}^2-mv_{1x}^2}{2}\,.$$
\vskip2.181bp\relax
\noindent (мы воспользовались формулой для $s_x$, выведенной в разделе «Равноускоренное движение»).

Заметим теперь, что в данном случае проекция скорости отличается от модуля скорости разве что знаком; поэтому $v_{1x}^2=v_1^2$ и $v_{2x}^2=v_2^2$. В результате имеем:
\vskip-2.833bp\relax
$$A=\frac{mv_2^2}{2}-\frac{mv_1^2}{2}=K_2-K_1=\Delta K,$$
\vskip2.181bp\relax
\noindent что и требовалось.

На самом деле соотношение $\Delta K=A$ справедливо и в самом общем случае криволинейного движе\-ния под действием переменной силы.

\vskip5.146bp\relax
\noindent\textbf{Теорема о кинетической энер\-гии.} Изменение кинетической энер\-гии тела равно работе, со\-вершённой приложенными к телу внешними силами за рассматриваемый промежуток времени.
\newpage

% Страницы 83--86. Текст и разметка восстановлены по предоставленному PDF.
Если работа внешних сил положительна, то кинетическая энергия увеличивается ($\Delta K>0$, тело разгоняется).

Если работа внешних сил отрицательна, то кинетическая энергия уменьшается ($\Delta K<0$, тело замедляет движение). Пример~--- торможение под действием силы трения, работа ко\-торой отрицательна.

Если же работа внешних сил равна нулю, то кинетическая энергия тела за это время не ме\-няется. Нетривиальный пример~--- равномерное движение по окружности, совершаемое грузом на нити в горизонтальной плоскости. Сила тяже\-сти, сила реакции опоры и сила натяжения нити всегда перпендикулярны скорости, и работа каждой из этих сил равна нулю в течение любого промежутка времени. Соответственно, кинетическая энергия груза (а значит, и его скорость) остаётся постоянной в процессе движе\-ния.

\vspace{10.331pt}\noindent
\textbf{Задача.} Автомобиль едет по горизонтальной дороге со скоростью $v$ и начинает резко тормо\-зить. Найти путь $s$, пройденный автомобилем до полной остановки, если коэффициент трения шин о дорогу равен $\mu$.

\vspace{5.165pt}\noindent
\textit{Решение.} Начальная кинетическая энергия автомобиля $K_1=\dfrac{mv^2}{2}$, конечная кинетическая энергия $K_2=0$. Изменение кинетической энер\-гии $\Delta K=K_2-K_1=-\dfrac{mv^2}{2}$.

На автомобиль действуют сила тяже\-сти $m\vec g$, реакция опоры $\vec N$ и сила трения $\vec f$. Сила тяже\-сти и реакция опоры, будучи перпендикулярны перемещению автомобиля, работы не соверша\-ют. Работа силы трения:
\[
A=-fs=-\mu Ns=-\mu mgs.
\]

Из теоремы о кинетической энер\-гии теперь получаем:
\[
\Delta K=A\Rightarrow-\frac{mv^2}{2}=-\mu mgs\Rightarrow s=\frac{v^2}{2\mu g}\,.
\]
\subsection*{1.16.5\quadПотенциальная энергия тела вблизи поверхности Земли}
\noindent Рассмотрим тело массы $m$, находящееся на некоторой высоте над поверхностью Земли. Высоту считаем много меньше земного радиуса. Изменением силы тяже\-сти в процессе перемещения тела пренебрегаем.

\begin{wrapfigure}{r}{189.306bp}
\vspace{-2.501pt}
\centering\hypertarget{fig150}{}\ReplicaGravityDiagram
\par\vspace{10.05pt}\hspace*{3.916bp}Рис.~1.50. $A=mg(h_1-h_2)$
\vspace{-6pt}
\end{wrapfigure}
Если тело находится на высоте $h$, то \textit{потенциаль\-ная энергия} тела по определению равна:
\[
W=mgh,
\]
\noindent где $g$~--- ускорение свободного падения вблизи поверхности Земли.

Высоту не обязательно отсчитывать от поверхности Земли. Как мы увидим ниже (формулы (\hyperlink{eq178}{1.78}), (\hyperlink{eq179}{1.79})), фи\-зическим смыслом обладает не сама по себе потенциаль\-ная энергия, но её \textit{изменение}. А изменение потенциаль\-ной энер\-гии не зависит от уровня отсчёта. Выбор нуле\-вого уровня потенциаль\-ной энер\-гии в конкретной задаче диктуется исключительно соображениями удобства.

Найдём работу, совершаемую силой тяже\-сти при перемещении тела. Предположим, что тело перемещается по прямой из точки $P$, находящейся на высоте $h_1$, в точку $Q$, находящуюся на высоте $h_2$ (рис.~\hyperlink{fig150}{1.50}).
\newpage
Угол между силой тяже\-сти $m\vec g$ и перемещением тела $\vec s$ обозначим $\alpha$. Для работы силы тяже\-сти получим:
\[
A=m\vec g\vec s=mgs\cos\alpha.
\]
\noindent Но, как видно из рис.~\hyperlink{fig150}{1.50}, $s\cos\alpha=h_1-h_2$. Поэтому
\[
A=mg(h_1-h_2)=mgh_1-mgh_2,
\]
\noindent или
\begin{equation}\hypertarget{eq178}{}A=W_1-W_2.\tag{1.78}\end{equation}
\noindent Учитывая, что $W_1-W_2=-(W_2-W_1)=-\Delta W$, имеем также:
\begin{equation}\hypertarget{eq179}{}A=-\Delta W.\tag{1.79}\end{equation}

Можно доказать, что формулы (\hyperlink{eq178}{1.78}) и (\hyperlink{eq179}{1.79}) справедливы \textit{для любой тра\-ек\-тории}, по ко\-торой тело перемещается из точки $P$ в точку $Q$, а не только для прямолинейного отрезка.

\vspace{5.165pt}
\textit{Работа силы тяже\-сти не зависит от формы тра\-ек\-тории, по ко\-торой перемещается те\-ло, и равна разности значений потенциаль\-ной энер\-гии в начальной и конечной точках тра\-ек\-тории. Иными словами, работа силы тяже\-сти всегда равна изменению потенциаль\-ной энер\-гии с противоположным знаком.}

\textit{В частности, работа силы тяже\-сти по любому замкнутому пути равна нулю.}

\vspace{5.165pt}
Сила называется \textit{консервативной}, если при перемещении тела работа этой силы не зависит от формы тра\-ек\-тории, а определяется только начальным и конечным положением тела. Сила тяже\-сти, таким образом, является консервативной. Работа консервативной силы по любому замкнутому пути равна нулю. Только в случае консервативной силы возможно ввести такую величину, как потенциаль\-ная энергия.

\subsection*{1.16.6\quadПотенциальная энергия деформированной пружины}
\noindent Рассмотрим пружину жёсткости $k$. Начальная деформация пружины равна $x_1$. Предположим, что пружина деформируется до некоторой конечной величины деформации $x_2$. Чему равна при этом работа силы упру\-гости пружины?

В данном случае силу на перемещение не умножишь, так как сила упру\-гости ме\-няется в процессе деформации пружины. Для нахождения работы переменной силы требуется интегри\-рование. Мы не будем приводить здесь вывод, а сразу выпишем конечный результат.

Оказывается, сила упру\-гости пружины также является консервативной. Её работа зависит лишь от величин $x_1$, $x_2$ и определяется формулой:
\[
A=\frac{kx_1^2}{2}-\frac{kx_2^2}{2}\,.
\]

Величина
\[
W=\frac{kx^2}{2}
\]
\noindent называется \textit{потенциаль\-ной энергией деформированной пружины} ($x$~--- величина деформации). Следовательно,
\[
A=W_1-W_2=-\Delta W,
\]
\noindent что полностью аналогично формулам (\hyperlink{eq178}{1.78}) и (\hyperlink{eq179}{1.79}).
\newpage
\subsection*{1.16.7\quadЗакон сохране\-ния механи\-ческой энер\-гии}
\noindent Консервативные силы называются так потому, что сохраняют механическую энергию замкну\-той си\-стемы тел.

\textit{Механическая энергия} $E$ тела равна сумме его кинетической и потенциаль\-ной энергий:
\[
E=K+W.
\]

Механическая энергия си\-стемы тел равна сумме их кинетических энергий и потенциаль\-ной энер\-гии их взаимодействия друг с другом.

Предположим, что тело совершает движение под действием силы тяже\-сти и/или силы упру\-гости пружины. Будем считать, что трения нет. Пусть в начальном положении кинетическая и потенциаль\-ная энер\-гии тела равны $K_1$ и $W_1$, в конечном положении~--- $K_2$ и $W_2$. Работу внешних сил при перемещении тела из начального положения в конечное обозначим $A$.

По теореме о кинетической энер\-гии:
\[
K_2-K_1=A.
\]

Но работа консервативных сил равна разности потенциальных энергий:
\[
A=W_1-W_2.
\]

Отсюда получаем:
\[
K_2-K_1=W_1-W_2,
\]
\noindent или
\[
K_1+W_1=K_2+W_2.
\]

Левая и правая части данного равенства представляют собой механическую энергию тела в начальном и конечном положении:
\[
E_1=E_2.
\]

Следовательно, \textit{при движении тела в поле силы тяже\-сти и/или на пружине механиче\-ская энергия тела остаётся неизменной при отсутствии трения.}

Справедливо и более общее утверждение.

\vspace{5.165pt}\noindent
\textbf{Закон сохране\-ния механи\-ческой энер\-гии.} Если в замкну\-той системе действуют только консервативные силы, то механиче\-ская энергия си\-стемы сохраняется.

\vspace{5.165pt}
При этих условиях могут происходить лишь превращения энер\-гии: из кинетической в по\-тенциальную и наоборот. Общий запас механи\-ческой энер\-гии си\-стемы остаётся постоянным.

\subsection*{1.16.8\quadЗакон изменения механи\-ческой энер\-гии}
\noindent Если между телами замкну\-той си\-стемы имеются силы сопротивления (сухое или вязкое тре\-ние), то механиче\-ская энергия си\-стемы будет уменьшаться. Так, автомобиль останавливается в результате торможения, колебания маятника постепенно затухают и т.\,д. Силы трения неконсер\-вативны: работа силы трения очевидным образом зависит от пути, по которому перемещается тело между данными точками. В частности, работа силы трения по замкнутому пути не равна нулю.

Снова рассмотрим движение тела в поле силы тяже\-сти и/или на пружине. Вдобавок на тело действует сила трения, которая за рассматриваемый промежуток времени совершает отрица\-тельную работу $A_{\text{тр}}$. Работу консервативных сил (тяже\-сти и упру\-гости) по-прежнему обозна\-чаем $A$.
\newpage
Изменение кинетической энер\-гии тела равно работе \textit{всех} внешних сил:
\[
K_2-K_1=A+A_{\text{тр}}.
\]

Но $A=W_1-W_2$, следовательно
\[
K_2-K_1=W_1-W_2+A_{\text{тр}}.
\]

Отсюда
\[
K_2+W_2-(K_1+W_1)=A_{\text{тр}},
\]
\noindent или
\[
E_2-E_1=A_{\text{тр}}.
\]

В левой части стоит величина $\Delta E=E_2-E_1$~--- изменение механи\-ческой энер\-гии тела:
\[
\Delta E=A_{\text{тр}}.
\]

Итак, \textit{при движении тела в поле силы тяже\-сти и/или на пружине изменение механи\-ческой энер\-гии тела равно работе силы трения.} Так как работа силы трения отрицательна, изменение механи\-ческой энер\-гии также отрицательно: механиче\-ская энергия убывает.

Справедливо и более общее утверждение.

\vspace{5.165pt}\noindent
\textbf{Закон изменения механи\-ческой энер\-гии.} Изменение механи\-ческой энер\-гии замкну\-той си\-стемы равно работе сил трения, действующих внутри си\-стемы.

\vspace{5.165pt}
Ясно, что закон сохране\-ния механи\-ческой энер\-гии является частным случаем данного утвер\-ждения.

Конечно, убыль механи\-ческой энер\-гии не противоречит общефизическому закону сохране\-ния энер\-гии. В данном случае механиче\-ская энергия превращается в энергию теплового движе\-ния частиц вещества и их по\-тенциальную энергию взаимодействия друг с другом, т.~е. переходит во \textit{внутреннюю энергию} тел си\-стемы.

\end{document}
